\section{A propositional countervaluation method}
\label{sec:formal-practice-02}

To refute \(\Gamma\entails\varphi\), the target conditions are
\[
  \widehat v(\gamma)=\True\quad(\gamma\in\Gamma),
  \qquad
  \widehat v(\varphi)=\False.
\]
Work inward from these target values.  A false implication requires a true
antecedent and false consequent; a true conjunction requires two true
conjuncts; a false disjunction requires two false disjuncts.  When a branch
assigns both truth values to the same propositional variable, close that branch.
A surviving total valuation is a countervaluation.

\begin{example}
For \((P\to Q)\lor(Q\to P)\), an attempted false valuation requires both
implications false.  The first requires \(P\) true and \(Q\) false; the second
requires \(Q\) true and \(P\) false.  The requirements conflict, so no
countervaluation exists.  This search establishes validity only because every
way for the outer disjunction to be false was exhausted.
\end{example}

